\section{Article Search and Retrieval}\label{app:search_queries}

This section details the search query construction, database-specific adaptations, and PDF retrieval strategy used for article acquisition across priority pathogens in the \name~pipeline. Following the Pathogen Epidemiology Review Group (PERG) methodology\footnote{\url{https://github.com/mrc-ide/priority-pathogens/wiki/Search-terms}}, we developed a standardised base query structure that captures core epidemiological domains including transmission dynamics, disease severity, temporal parameters, transmission heterogeneity, and evolutionary characteristics.

Different bibliographic databases support different search capabilities, requiring tailored query implementations. We maintain two versions of each pathogen query: one for PubMed and Europe PMC (which support wildcard truncation operators using \texttt{*}), and another for OpenAlex (which requires fully expanded term variants).

\subsection{Base Search Query (PubMed and Europe PMC)}

The base query for PubMed and Europe PMC uses Boolean operators with truncation symbols to capture morphological term variations:

\begin{lstlisting}[basicstyle=\small\ttfamily, breaklines=true, backgroundcolor=\color{gray!10}]
[PATHOGEN_IDENTIFIER] AND (
    (transmissi* OR epidemiolog*) OR 
    (model* NOT imag*) OR 
    (severity OR "case fatality ratio*" OR CFR OR "case fatality rate*" 
        OR "mortality rate*" OR "attack rate*") OR 
    ("infectious period*" OR "serial interval*" OR "incubation period*" 
        OR "generation time*" OR "generation interval*" OR "latent period*" 
        OR latency) OR 
    (heterogeneit* OR superspread* OR "super spread*" OR super-spread* 
        OR overdispersion OR overdispersed OR over-dispersion OR over-dispersed 
        OR "over dispersion" OR "over dispersed") OR 
    (infectivity OR infectiousness OR "growth rate*" OR "reproduction number*" 
        OR "reproductive number*" OR R0 OR "reproduction ratio*" 
        OR "reproductive rate*") OR 
    ("pre-existing immunity" OR serological OR serology OR serosurvey*) OR 
    (evolution* OR mutation* OR substitution*) OR 
    (outbreak* OR cluster* OR epidemic*) OR 
    ("risk factor*")
    [ADDITIONAL_TERMS]
) [EXCLUSION_CRITERIA]
\end{lstlisting}

\subsection{OpenAlex Adapted Queries}

Because the OpenAlex API does not support wildcard operators\footnote{\url{https://docs.openalex.org/how-to-use-the-api/get-lists-of-entities/search-entities}} and strips these characters during query processing, we expanded all truncated terms into their common morphological variants:

\begin{lstlisting}[basicstyle=\small\ttfamily, breaklines=true, backgroundcolor=\color{gray!10}]
[PATHOGEN_IDENTIFIER] AND (
    (transmission OR transmissibility OR transmissible OR transmitted 
        OR transmitting OR transmit OR epidemiology OR epidemiological 
        OR epidemiologic) OR 
    (model OR models OR modeling OR modelling OR modeled OR modelled 
        NOT (image OR images OR imaging)) OR 
    (severity OR "case fatality ratio" OR "case fatality ratios" OR CFR 
        OR "case fatality rate" OR "case fatality rates" OR "mortality rate" 
        OR "mortality rates" OR "attack rate" OR "attack rates") OR 
    ("infectious period" OR "infectious periods" OR "serial interval" 
        OR "serial intervals" OR "incubation period" OR "incubation periods" 
        OR "generation time" OR "generation interval" OR "generation intervals" 
        OR "latent period" OR "latent periods" OR latency) OR 
    (heterogeneity OR heterogeneous OR superspread OR superspreader 
        OR superspreaders OR superspreading OR "super spread" 
        OR "super spreader" OR "super spreaders" OR "super spreading" 
        OR overdispersion OR overdispersed OR "over dispersion" 
        OR "over dispersed") OR 
    (infectivity OR infectiousness OR "growth rate" OR "growth rates" 
        OR "reproduction number" OR "reproduction numbers" 
        OR "reproductive number" OR "reproductive numbers" OR R0 
        OR "reproduction ratio" OR "reproduction ratios" 
        OR "reproductive rate" OR "reproductive rates" 
        OR "basic reproduction number") OR 
    ("pre-existing immunity" OR serological OR serology OR serosurvey 
        OR serosurveys OR seroprevalence OR serosurveillance) OR 
    (evolution OR evolutionary OR evolving OR evolved OR mutation 
        OR mutations OR mutant OR mutants OR mutate OR mutated 
        OR substitution OR substitutions) OR 
    (outbreak OR outbreaks OR cluster OR clusters OR clustering 
        OR epidemic OR epidemics OR pandemic OR pandemics) OR 
    ("risk factor" OR "risk factors")
    [ADDITIONAL_TERMS]
) [EXCLUSION_CRITERIA]
\end{lstlisting}

\subsection{Pathogen-Specific Query Modifications}

Table~\ref{tab:search_queries} summarises the pathogen-specific modifications applied across all database implementations. Most pathogens require only customised identifiers to ensure relevant literature retrieval. However, the queries for SARS explicitly exclude COVID-19 literature to prevent cross-contamination with SARS-CoV-2 studies. Similarly, queries for Zika include vector-specific epidemiological parameters (extrinsic incubation period, vector competence) that are essential for capturing mosquito-borne transmission dynamics. For Rift Valley fever, Crimean-Congo hemorrhagic fever (CCHF) and MERS, we incorporated additional virus-specific identifiers and spelling variants to enhance retrieval comprehensiveness. Despite these modifications, all databases share consistent pathogen identifiers and exclusion criteria, differing only in their use of wildcard forms (PubMed/Europe PMC) versus expanded term variants (OpenAlex).

\begin{table*}[h]
\footnotesize
\caption{\textbf{Pathogen-specific modifications to the standardised search query}. All databases share consistent pathogen identifiers and exclusion criteria; PubMed/Europe PMC use wildcard forms while OpenAlex uses expanded variants.}
\label{tab:search_queries}
\begin{center}
\begin{tabular}{lp{5cm}p{3.2cm}p{3.5cm}}
\toprule
\textbf{Pathogen} & \textbf{\texttt{PATHOGEN\_IDENTIFIER}} & \textbf{\texttt{ADDITIONAL\_TERMS}} & \textbf{\texttt{EXCLUSION\_CRITERIA}} \\
\midrule
Marburg virus & Marburg virus & --- & --- \\
Ebola virus & Ebola & --- & --- \\
Lassa virus & Lassa & --- & --- \\
SARS-CoV-1 & SARS OR SARS-CoV-1 OR ``Severe acute respiratory syndrome" & --- & NOT (COVID-19 OR SARS-CoV-2) \\
Zika virus & zika & OR (``extrinsic incubation period" OR ``EIP" OR ``vector competence" OR ``vectorial capacity")\textsuperscript{†} & --- \\
Nipah virus & Nipah & --- & --- \\
MERS-CoV & MERS OR MERS-CoV OR ``Middle East respiratory syndrome" OR ``Middle East Respiratory Syndrome Coronavirus"\textsuperscript{‡} & --- & --- \\
Rift Valley fever virus & ``Rift valley fever" OR RVF OR ``Rift Valley Fever Virus" OR RVFV\textsuperscript{‡} & --- & --- \\
CCHF virus & ``Crimean Congo haemorrhagic fever" OR ``Crimean-Congo hemorrhagic fever" OR CCHF OR ``CCHF virus" OR CCHFV\textsuperscript{‡} & --- & --- \\
\bottomrule
\end{tabular}
\end{center}
\vspace{-2mm}
{\small \textsuperscript{†}Vector-specific terms capture mosquito transmission parameters unique to arboviral epidemiology.}\\
{\small \textsuperscript{‡}Expanded identifiers include alternative spellings (American/British English), virus-specific nomenclature, and common abbreviations for comprehensive coverage.}
\end{table*}

\subsection{Metadata Extraction and Deduplication}

We extract bibliographic metadata from each database as summarised in Table~\ref{tab:metadata_fields}. OpenAlex provides direct PDF URLs and internal work identifiers, PubMed supplies standardised medical literature identifiers (PMID: PubMed ID; PMCID: PubMed Central ID), and Europe PMC offers full-text availability metadata. The Digital Object Identifier (DOI) serves as a persistent identifier across databases.

We implement a hierarchical five-level deduplication strategy:

\begin{enumerate}
\item \textbf{DOI-based}: Normalised DOI strings (case-insensitive, URL prefixes stripped);
\item \textbf{PMID-based}: Numeric PMID extraction and normalisation;
\item \textbf{PMCID-based}: Normalised PMC identifiers (uppercase, ``PMC'' prefix standardised);
\item \textbf{OpenAlex ID-based}: Internal OpenAlex work identifiers;
\item \textbf{Title-year combination}: Normalised title strings (lowercase, alphanumeric only) paired with publication year.
\end{enumerate}

When duplicate records are detected, identifier fields (DOI, PMID, PMCID, OpenAlex ID, URLs) preserve all non-null values while narrative fields (title, abstract, journal) retain the first non-null value. Source provenance is marked as ``Both" when records appear in multiple databases.

\begin{table}[h]
\centering
\footnotesize
\begin{minipage}{0.48\textwidth}
    \centering
    \caption{\textbf{Metadata fields extracted during article search.} PMID: PubMed ID; PMCID: PubMed Central ID; DOI: Digital Object Identifier.}
    \label{tab:metadata_fields}
    \begin{footnotesize}
    \begin{tabular}{ll}
    \toprule
    \textbf{Field} & \textbf{Description} \\
    \midrule
    article\_id & Generated unique identifier \\
    source & Database origin \\
    pmid & PubMed Identifier \\
    pmcid & PubMed Central Identifier \\
    doi & Digital Object Identifier \\
    title & Article title \\
    authors & Semicolon-delimited author list \\
    journal & Publication venue \\
    year & Publication year \\
    abstract & Article abstract \\
    url & Canonical article URL \\
    openalex\_id & OpenAlex work identifier \\
    openalex\_pdf\_url & Direct PDF link from OpenAlex \\
    pathogen & Target pathogen \\
    query & Search query used \\
    harvested\_at & ISO 8601 timestamp \\
    \bottomrule
    \end{tabular}
    \end{footnotesize}
\end{minipage}
\hfill
\begin{minipage}{0.48\textwidth}
    \centering
    \caption{\textbf{Additional fields populated during PDF retrieval attempts.}}
    \label{tab:download_fields}
    \begin{footnotesize}
    \begin{tabular}{ll}
    \toprule
    \textbf{Field} & \textbf{Description} \\
    \midrule
    downloaded & Boolean success flag \\
    downloaded\_path & Filesystem path to PDF \\
    download\_source & Source that provided PDF \\
    download\_attempted\_at & ISO 8601 timestamp \\
    download\_error & Error messages from attempts \\
    \bottomrule
    \end{tabular}
    \end{footnotesize}
\end{minipage}
\end{table}

\subsection{PDF Retrieval}\label{app:article_download}

We attempt PDF downloads through multiple open access sources using a cascading retrieval strategy. Before attempting downloads, available identifiers (PMID, PMCID, DOI) are cross-referenced using NCBI's PMC ID Converter API\footnote{\url{https://www.ncbi.nlm.nih.gov/pmc/tools/id-converter-api/}} to maximise source compatibility. The system then attempts downloads from up to four sources in priority order (Table~\ref{tab:download_sources}), stopping at the first successful retrieval.

\subsubsection{Implementation Details}

Downloads employ HTTP streaming to temporary files with 64~KB chunks and validate each file through two stages: (1) magic byte verification (\texttt{\%PDF} header), and (2) content inspection for HTML access denial pages. Files exceeding 500~MB or failing validation are immediately discarded. Thread-pool parallelism with 16 workers processes downloads concurrently while respecting per-source rate limits. In-memory caches keyed by normalised identifiers store both successful PDF URLs and negative markers to eliminate redundant API calls. Progress is checkpointed every 50 records for crash recovery.

Successfully validated PDFs are saved with standardised filenames following identifier priority (PMID $\rightarrow$ PMCID $\rightarrow$ DOI hash $\rightarrow$ title hash). Metadata is augmented with download provenance including source, timestamp, and error diagnostics.

\begin{table}[t]
\centering
\caption{\textbf{PDF retrieval sources in cascading priority order.} Sources are queried sequentially until success or exhaustion. Identifier cross-referencing via NCBI PMC ID Converter API precedes all download attempts (10 req/s, cached).}
\label{tab:download_sources}
\begin{small}
\begin{tabular}{clcc}
\toprule
\textbf{Priority} & \textbf{Source \& Endpoint} & \textbf{Rate Limit} & \textbf{Cached} \\
\midrule
1 & \textbf{OpenAlex Direct PDF URL} & 30 req/s & No \\
  & \quad Metadata field \texttt{openalex\_pdf\_url} & & \\
\midrule
2 & \textbf{Europe PMC Fulltext API} & 20 req/s & Yes \\
  & \quad \texttt{ebi.ac.uk/europepmc/webservices/rest/search} & & \\
\midrule
3 & \textbf{Unpaywall API} & 50 req/s & Yes \\
  & \quad \texttt{api.unpaywall.org/v2/\{DOI\}?email=\{EMAIL\}} & & \\
\midrule
4 & \textbf{OpenAlex DOI Lookup} & 30 req/s & Yes \\
  & \quad \texttt{api.openalex.org/works/https://doi.org/\{DOI\}} & & \\
\bottomrule
\end{tabular}
\end{small}
\end{table}

\subsection{Final Quality Control}

After retrieval, we applied deduplication and quality filtering that removes: records lacking abstracts, duplicate article IDs, duplicate DOIs (retaining first occurrence) and records with file validation failures.
